% !TeX root = main.tex
\section{Duty-Cycle Modulation via Phase Warping}
\subsection{Downstroke Ratio}
We define the downstroke ratio $\delta \in (0,1)$ as the fraction of a period $T$ spent on the downstroke half-stroke.
For a fixed effective frequency $f_{\mathrm{eff}}=1/T$, the corresponding half-stroke frequencies are
\begin{equation}
f_{\mathrm{down}} = \frac{f_{\mathrm{eff}}}{2\delta}, \qquad
f_{\mathrm{up}} = \frac{f_{\mathrm{eff}}}{2(1-\delta)}.
\end{equation}
When $\delta = 0.5$, the cycle is symmetric in time.

\subsection{Phase-Warped SCCPFM Waveform}
Let $u = (t \bmod T)/T \in [0,1)$ denote the normalized phase time within a cycle.
We define a phase map $p(u)$ that allocates $p \in [0,0.5]$ to the downstroke segment and $p \in [0.5,1]$ to the upstroke segment:
\begin{equation}
p(u) =
\begin{cases}
\frac{1}{2\delta}u, & u \le \delta,\\
\frac{1}{2} + \frac{1}{2(1-\delta)}(u-\delta), & u > \delta.
\end{cases}
\end{equation}
The warped phase is $\psi(t) = 2\pi p(u)$.
The commanded stroke angle is then defined using a cosine waveform
\begin{equation}
q(t) = A \cos(\psi(t)),
\end{equation}
with amplitude $A$.
This formulation makes $\delta$ a single scalar input that reshapes instantaneous velocity and acceleration profiles while preserving the stroke amplitude.

\subsection{Interpretation as Time Reparameterization}
Duty-cycle modulation does not alter the geometric trajectory in the $q$--$\psi$ domain; it alters how quickly the trajectory is traversed within different segments of the cycle.
Since aerodynamic forces are nonlinear in instantaneous kinematics and depend on phase-dependent effective angle of attack, time reparameterization can change cycle-averaged force and power even when the stroke amplitude is fixed.
